% =====================================================================
%  BIBLIOGRAFÍA (formato APA, transcrita del documento Word)
%  Cada referencia es un \item dentro del entorno referenciasapa.
% =====================================================================
\chapter*{BIBLIOGRAFÍA}
\addcontentsline{toc}{chapter}{BIBLIOGRAFÍA}
\label{cap:bibliografia}

\begin{referenciasapa}

\item Abrahamsson, P., Hanhineva, A., Hulkko, H., Ihme, T., Jäälinoja, J.,
Koskela, J., \ldots{} Salo, O. (2004). \textit{Mobile-D: An Agile Approach for
Mobile Application Development.} Helsinki: VTT Publications.

\item Alliance, O. H. (2022). Obtenido de Android Open Source Project
Overview.: \url{https://source.android.com/?hl=es-419}

\item Android Developers, G. (2024). \textit{Introducción a Android Studio}.
Obtenido de Android Developers:
\url{https://developer.android.com/studio/intro?hl=es-419}

\item Apple. (2023). \textit{Apple Developer Documentation}. Obtenido de
\url{https://developer.apple.com/}

\item AppMaster. (4 de septiembre de 2023). \textit{Programación Móvil}.
Obtenido de \url{https://appmaster.io/es/glossary/programacion-movil}

\item Apps, L. B. (2011). \textit{Guía de aplicaciones móviles.} Obtenido de
Guía de aplicaciones móviles. Fundación Telefónica.

\item Atlassian. (2023). \textit{About us: Trello history, logos \& customers
| Trello}. Obtenido de Atlassian Trello: \url{https://trello.com/about}

\item Azuma, R. (1997). \textit{A Survey of Augmented Reality.} Obtenido de
\url{https://www.cs.unc.edu/~azuma/ARpresence.pdf}

\item B, G. (13 de marzo de 2025). \textit{Qué es GitHub y cómo empezar a
usarlo}. Obtenido de ES Tutoriales:
\url{https://www.hostinger.com/es/tutoriales/que-es-github}

\item Beck, K., \& Andres, C. (2005). Extreme Programming Explained: Embrace
change. En \textit{Extreme Programming Explained.} Addison-Wesley.

\item Booch, G., Maksimchuk, R., Engel, M., Young, B., Conallen, J., \&
Houston, K. (2007). \textit{Object-Oriented Analysis and Design with
Applications.} Addison-Wesley.

\item Cloud, G. (2025). \textit{Inteligencia Artificial (IA): una guía fácil de
entender}. Obtenido de
\url{https://cloud.google.com/learn/what-is-artificial-intelligence?hl=es-419}

\item Corporation, O. (2024). \textit{MySQL Documentation}. Obtenido de MySQL:
\url{https://dev.mysql.com/doc/}

\item Cuadrado, G. C. (22 de julio de 2022). \textit{OpenWebinars}. Obtenido de
\url{https://openwebinars.net/blog/que-es-visual-studio-code-y-que-ventajas-ofrece/}

\item \textit{Curso para desarrollo de apps con React Native}. (s.f.). Obtenido
de Código Facilito: \url{https://codigofacilito.com/articulos/que-es-react-native}

\item Deloitte. (2019). \textit{¿Qué es React Native? | Tecnología | Deloitte
España}. Obtenido de Deloitte:
\url{https://www.deloitte.com/es/es/services/consulting/blogs/todo-tecnologia/que-es-react-native.html}

\item Developers, G. F. (2025). \textit{Crea nuevas experiencias de realidad
aumentada que combinen perfectamente con el mundo digital con el físico |
ARCore | Google for Developers.} Obtenido de ARCore:
\url{https://developers.google.com/ar?hl=es-419}

\item DuBois, P. (2013). \textit{MySQL.} Addison-Wesley Professional.

\item Edward Laux, J. S. (2016). Obtenido de BodBot AI Entrenador Personal:
\url{https://www.bodbot.com/}

\item Flutter, s. (s.f.). \textit{Flutter - Crea aplicaciones para cualquier
pantalla.} Obtenido de \url{https://esflutter.dev/}

\item Gabriel Maide, E., \& Pacienzia, J. (noviembre de 2020).
\textit{Metodologías de desarrollo de software aplicadas a proyectos de
automatización de procesos.} Obtenido de Pistas Educativas:
\url{https://pistaseducativas.celaya.tecnm.mx/index.php/pistas/article/viewFile/2381/1900}

\item GeeksforGeeks. (11 de julio de 2025). \textit{Incremental Process Model
Software Engineering.} Obtenido de GeeksforGeeks:
\url{https://www.geeksforgeeks.org/software-engineering/software-engineering-incremental-process-model/}

\item Gironés, T. (2012). \textit{Aplicaciones móviles: Conceptos y
desarrollo.}

\item Git. (s.f.). Obtenido de \url{http://git-scm.com/}

\item Gjestvang, C. S. (s.f.). \textit{Motives and barriers to initiation and
sustained exercise adherence in new fitness club members: A prospective study.
Publicado en Frontiers in Psychology.} Obtenido de
\url{https://psycnet.apa.org/record/2020-64605-023}

\item Global, S. (noviembre de 2025). Obtenido de Mobile Operating System Market
Share Plurinational State of Bolivia | StatCounter Global Stats:
\url{https://gs.statcounter.com/os-market-share/mobile/bolivia}

\item Gómez, \& Hernández. (2016). Mobile-D: Metodología ágil para el desarrollo
de aplicaciones móviles. \textit{Revista Orbis}.

\item Gómez, A. A. (2015). \textit{Comparación de dos tecnologías de desarrollo
de aplicaciones móviles desde la perspectiva de los atributos de calidad}.
Obtenido de ResearchGate:
\url{https://www.researchgate.net/figure/Figura-5-Arquitectura-de-Android_fig3_319975215}

\item Gómez, A. A. (2015). \textit{Comparación de dos tecnologías de desarrollo
de aplicaciones móviles desde la perspectiva de los atributos de calidad.}
Universidad Nacional de Colombia: Bogotá.

\item Google. (2023). \textit{Android Developers -- Official Documentation}.
Obtenido de \url{https://developer.android.com/}

\item Group, L. F. (2021). \textit{Ejercicios en el Gimnasio}. Obtenido de
\url{https://play.google.com/store/apps/details?id=gymworkout.gym.gymlog.gymtrainer&hl=es_BO}

\item Guide, I. S. (12 de septiembre de 2023). \textit{Modelo en espiral:
modelo de reducción de riesgos para proyectos de software.} Obtenido de
\url{https://www.ionos.mx/startupguide/productividad/modelo-en-espiral/}

\item Gurnov, A., \& Wrike (W.). (13 de noviembre de 2025). \textit{Las 21
mejores herramientas de gestión de proyectos}. Obtenido de Wrike:
\url{https://www.wrike.com/es/project-management-guide/faq/que-son-las-herramientas-de-gestion-de-proyectos/}

\item Hayes, M., \& Downie, A. (28 de noviembre de 2025). \textit{Realidad
aumentada. IBM}. Obtenido de IBM:
\url{https://www.ibm.com/es-es/think/topics/augmented-reality}

\item Huang, S. (2020). \textit{NoSQL Databases: Concepts and Applications.}
Cambridge: Morgan Kaufmann.

\item IBM. (2025). \textit{¿Qué son las redes neuronales?} Obtenido de
\url{https://www.ibm.com/es-es/think/topics/neural-networks}

\item IEEE. (1990). \textit{IEEE Standard Glossary of Software Engineering
Terminology}. Obtenido de Institute of Electrical and Electronics Engineers.

\item Inc., A. (2025). \textit{ARKit 6 - Augmented Reality - Apple Developer.}

\item Ionic Documentation, F. (2023). \textit{Ionic Framework.} Obtenido de
Ionic Documentation: \url{https://ionicframework.com/docs}

\item Isaac. (2025). \textit{Firebase: Qué es y cómo funciona la plataforma de
Google}. Obtenido de Mundobytes: \url{https://mundobytes.com/firebase-que-es/}

\item Kolt, G., \& Snyder-Mackler, L. (2007). \textit{Terapias físicas en el
deporte y el ejercicio.}

\item Leyva, R., León, G., \& García, F. (2016). Metodología Mobile-D para el
desarrollo de aplicaciones móviles. \textit{Revista Cubana de Ciencias
Informáticas, 10}(3), 45--57. Obtenido de
\url{https://www.researchgate.net/figure/Figura-210-Metodologia-Mobile-D-Fuente-Leyva-et-al-2016_fig5_348295603}

\item López Mora, S. (noviembre de 2025). \textit{Firebase: qué es, para qué
sirve, funcionalidades y ventajas}. Obtenido de DIGITAL55:
\url{https://digital55.com/blog/que-es-firebase-funcionalidades-ventajas-conclusiones/}

\item Meier, R., \& Lake, I. (2018). \textit{Professional Android} (4.ª ed.).

\item Microsoft. (2022). \textit{Xamarin Documentation}. Obtenido de Microsoft
Learn: \url{https://learn.microsoft.com/en-us/xamarin/}

\item Mitchell, T. (1997). \textit{Machine Learning.} New York: McGraw-Hill.

\item MongoDB, I. (2024). \textit{MongoDB Atlas documentation}. Obtenido de
\url{https://www.mongodb.com/docs/atlas/}

\item Navarro, F., Martínez, A., \& Martínez, J. (2018). \textit{Realidad
Virtual y Realidad Aumentada.} Ra-Ma.

\item Nielsen, J. (1993). \textit{Usability Engineering.} Academic Press.

\item Nystrom, K. (2021). \textit{Kotlin Programming: The Big Nerd Ranch Guide.
Big Nerd Ranch.} Pearson Education Store.

\item Orientanet. (s.f.). \textit{Orientanet}. Obtenido de ¿Qué es
entrenamiento físico y tipos?:
\url{https://www.orientanet.es/que-es-entrenamiento-fisico-y-tipos/}

\item Pressman, R. S. (2010). \textit{Software Engineering: A Practitioner's
Approach.} Higher Education.

\item Pykes, K. (2024). \textit{Aprendizaje automático en Python: Tutorial de
Scikit-Learn}. Obtenido de
\url{https://www.datacamp.com/es/tutorial/machine-learning-python}

\item \textit{React Native}. (2025). Obtenido de
\url{https://reactnative.dev/docs/intro-react-native-components}

\item Santander, B. (21 de diciembre de 2020). \textit{Metodologías de
desarrollo de software: ¿qué son?} Obtenido de Santander:
\url{https://www.santanderopenacademy.com/es/blog/metodologias-desarrollo-software.html}

\item Sanz, D. (2019). Obtenido de Introducción a Android:
\url{https://dspace.itsjapon.edu.ec/jspui/bitstream/123456789/434/1/introduccion-android.pdf}

\item Schwaber, K., \& Sutherland, J. (2020). \textit{The Scrum Guide: The
Definitive Guide to Scrum.} Boston: Scrum.org.

\item Services, A. W. (2025). Obtenido de ¿Qué es la IA? -- Explicación de la
inteligencia artificial -- AWS:
\url{https://aws.amazon.com/es/what-is/artificial-intelligence/}

\item Sharma, M. (2018). \textit{Productivity gains with modern IDEs: A review
of IntelliJ-based environments.}

\item Smith, J., \& Darrow, M. (2019). \textit{Gradle for Android Development.}
O'Reilly Media.

\item Sommerville, I. (2011). \textit{Software Engineering.} Boston:
Addison-Wesley.

\item SQLite Consortium. (2024). \textit{SQLite Home page}. Obtenido de
SQLite.org: \url{https://sqlite.org/}

\item Team. (2024). \textit{Metodologías de desarrollo de software: guía para
optimizar tu proyecto.} Obtenido de Blog de Desarrollo de Software, Testing e
Inteligencia Artificial | Abstracta:
\url{https://es.abstracta.us/blog/metodologias-de-desarrollo-de-software-guia-completa/}

\item Tomás, J. (2019). \textit{El gran libro de Android} (Vol. 7). Bogotá:
Alfaomega.

\item Viera, U. (s.f.). \textit{Introducción a React Native: Crea Apps Móviles
Nativas en iOS y Android}. Obtenido de urianviera.com:
\url{https://urianviera.com/react-native/introduccion-a-react-native}

\item Yépez, M. D. (2017). \textit{Club Ensayos}. Obtenido de
\url{https://www.clubensayos.com/Negocios/Proyecto-de-grado-Creaci\%C3\%B3n-de-un-centro-de/4574087.html}

\end{referenciasapa}
